\section{Frozen packet and semilocal normalization}
\label{sec:packet}

\subsection{Transform and polarization conventions}

The inner product on \(L^2(\R,dx)\) is antilinear in the first variable:
\[
\langle\Phi,\Psi\rangle
=\int_{\R}\overline{\Phi(x)}\Psi(x)\,dx.
\]
Our Fourier transform and inverse normalization are
\[
\widehat\Phi(z)=\int_{\R}\Phi(x)e^{-izx}\,dx,
\qquad
\langle\Phi,\Psi\rangle
=\frac{1}{2\pi}\int_{\R}
\overline{\widehat\Phi(t)}\widehat\Psi(t)\,dt.
\]
Translation \(T_y\Phi(x)=\Phi(x-y)\) obeys
\[
\widehat{T_y\Phi}(t)=e^{-iyt}\widehat\Phi(t).
\]
For multiplicative tests we use
\[
\widetilde f(s)=\int_0^\infty f(u)u^{s-1}\,du.
\]

For \(\Phi,\Psi\in C_c^\infty(\R)\), define
\[
C_{\Phi,\Psi}(r)
=\int_{\R}\overline{\Phi(x-r)}\Psi(x)\,dx.
\]

\subsection{Complete semilocal Weil preform}

The source normalization reconstructed from
\cite{ConnesConsani2023,CCM2025} is
\begin{align}
\qform(\Phi,\Psi)
={}&
\frac{1}{2\pi}\int_{\R}
m_\infty(t)\,
\overline{\widehat\Phi(t)}\widehat\Psi(t)\,dt
\notag\\
&+
\overline{\widehat\Phi(-i/2)}\widehat\Psi(i/2)
+
\overline{\widehat\Phi(i/2)}\widehat\Psi(-i/2)
\notag\\
&-
\sum_{\substack{p\ {\rm prime}\\m\ge1}}
\frac{\log p}{p^{m/2}}
\left\{
C_{\Phi,\Psi}(m\log p)
+
C_{\Phi,\Psi}(-m\log p)
\right\},
\label{eq:preform}
\end{align}
where
\[
m_\infty(t)
=
\Re\frac{\Gamma'}{\Gamma}
\left(\frac14+\frac{it}{2}\right)-\log\pi.
\]
Every prime power \(p^m\), rather than only every prime, appears with exact
coefficient
\[
-(\log p)p^{-m/2}
=-\Lambda(p^m)(p^m)^{-1/2}.
\]
For a finite semilocal interval \([-\log\lambda,\log\lambda]\), the cutoff is
\[
p^m\le\lambda^2.
\]
The endpoint is inclusive and carries no half weight.

The equivalent real-place distribution is
\begin{align}
W_{\R}(C)
={}&
(\log 4\pi+\gamma_E)C(0)
\notag\\
&+
\int_0^\infty
\frac{
e^{r/2}\{C(r)+C(-r)\}-2C(0)
}{
e^r-e^{-r}
}\,dr,
\label{eq:real-place}
\end{align}
and the archimedean contribution to \eqref{eq:preform} is
\(-W_{\R}\).  This fixes the scalar, subtraction, and sign.

\subsection{The dyadic infinite-convolution packet}

For \(a_n=2^{-n-1}\), \(n\ge1\), put
\[
b_{a_n}(x)=\frac{1}{2a_n}\mathbf1_{[-a_n,a_n]}(x)
\]
and define
\[
g=*_{n\ge1}b_{a_n},\qquad
h=\frac{g}{c_g},\qquad c_g=\lVert g\rVert_2.
\]
Since \(\sum_{n\ge1}a_n=1/2\), the infinite convolution is the density
of an absolutely convergent sum of independent compactly supported random
variables.  Moreover
\[
g\in C_c^\infty(\R),\qquad
\supp g=[-1/2,1/2],
\]
and \(g\) is real, even, nonnegative, and strictly positive in the
interior.  Smoothness follows from the super-polynomial Fourier decay
obtained by bounding any fixed number of sinc factors.

The transform is
\[
\widehat h(z)
=c_g^{-1}
\prod_{n\ge1}\frac{\sin(a_nz)}{a_nz}.
\]
The product converges locally uniformly, and
\[
Z(\widehat h)=4\pi\Z\setminus\{0\}.
\]

Because \(h\) is even,
\[
K_h(r)=\int_\R h(x-r)h(x)\,dx=(h*h)(r).
\]
It is real, even, nonnegative, smooth, supported exactly on \([-1,1]\),
positive on \((-1,1)\), and \(K_h(0)=1\).  Its bilateral Laplace transform
is
\begin{equation}
\Lh(w)
=c_g^{-2}
\prod_{n\ge1}
\left(\frac{\sinh(a_nw)}{a_nw}\right)^2.
\label{eq:L-product}
\end{equation}
This is entire and even.  Since
\[
\frac{\sinh(a_nw)}{a_nw}=1+O(a_n^2)
\]
uniformly on compact \(w\)-sets and \(\sum a_n^2<\infty\), the infinite
tail has no extra zeros.  Hence
\begin{equation}
Z(\Lh)=4\pi i\Z\setminus\{0\}.
\label{eq:L-zero-set}
\end{equation}
At \(4\pi i k\ne0\), the multiplicity is
\[
2\bigl(v_2(|k|)+1\bigr).
\]
